\chapterpage{diagnosis}{读懂强雨失效的三层证据}{先问指标在评估哪个输出，再决定改模型还是改读出}
\label{p-diagnosis}
\diagram[0.83]{04-diagnosis.pdf}{4}{发生、条件幅值、尾部识别要分别回答。第三条轨道还应区分节点越线、概率排序与固定点值检出。}

\textbf{发生层。} \term{brier}{Brier} 评估有雨概率与真实有雨指示的平方差。\term{rank}{AP、AUROC}关注分数排序。若正例定义为有雨，指标只回答有雨问题；AP20、AUROC20 则把强雨作为正例。\src{P2}\src{P8}\src{R6}

\textbf{条件幅值层。} \term{pinball}{Pinball} 衡量指定分位点的误差。REAL 相对 ZERO 在同一 seed、固定 E9、2024 Dev 的高海拔强雨子集上，物理 Pinball/20 降 0.1367\%，log Pinball 降 0.2106\%。该子集含 178 个\term{event}{像元时}，落在 80 个场景内。两个较低海拔层方向相反，高海拔 \(>50\,\mmh\) 样本为零。全域 log 损失略升，Brier 略降，AP 和 AUROC 略降。两臂三种点估计均命中 0/717，高海拔 q32 覆盖由 6/178 增至 11/178。该比较支持单 seed 的局部描述。\src{P6}

\textbf{尾部识别层。} \term{tail}{q32 覆盖与节点越线}各有自己的判断条件。TAIL20 原九项筛查通过八项。强雨物理 Pinball/20 从 0.506865 降到 0.504120，真强雨 q32 越线 160 增至 179。唯一失守项是非强雨 q32 越线，5,769 增至 6,602。A 点值仍命中 0/717，候选封存。\src{P7}

\sub{A20 的命中增加，要怎样解释？}
同一批 963 个场景、3,303,090 个有效像元时中，强雨正例 717 个。下表评价同一固定点值 A 的阈值检出和跨阈值排序。\src{P8}\par
\begin{tabularx}{\linewidth}{@{}Lrr@{}}
\toprule
指标 & COSINE9 & A20 \\
\midrule
固定 \(>20\,\mmh\) 的\term{confusion}{命中 / 误报} & 0 / 0 & 627 / 177,098 \\
全体 AP20 & 0.018691 & 0.007786 \\
全体 AUROC20 & 0.967601 & 0.968501 \\
海拔 \(\ge 2{,}000\) m AP20 & 0.008314 & 0.004929 \\
海拔 \(\ge 2{,}000\) m AUROC20 & 0.978222 & 0.971738 \\
\bottomrule
\end{tabularx}

A20 命中率约 87.45\%，报出的强雨中约 0.35\% 是真强雨。全体 AUROC20 略升，但两个已报告范围的 AP20 均下降。原筛查九项仅过两项，固定阈值误报线失守。排序补证后按预设规则结束该候选。点值排序在这里仍需和\term{calibration}{概率校准}分开识别。\src{P8}\src{P13}\par\hyperlink{case-figures}{打开两幅结果图 $\rightarrow$}

\clearpage
\hypertarget{case-figures}{}
\section*{用两幅结果图练习判读}
{\color{accent}\bfseries 先找比较器、样本量和基准线，再读差值方向。}\par

\diagram{05-terrain.pdf}{5}{REAL 相对 ZERO 的强雨分位损失变化。负值表示损失下降。三个海拔层的 \(n\) 都是像元时数。固定 B7/seed2029/E9 的 2024 开发配对，未给出独立天气过程数或跨种子区间。\src{P6}}
\textbf{读图动作。} 先按同一指标横向看三个海拔层，再比较同一层的两个损失尺度。选定高海拔子集后，句子里保留中低海拔的反例与高海拔极端样本不足的限制。\par

\diagram{06-ranking.pdf}{6}{同一强雨定义下的 A 点值 AP20。参考为 COSINE9，候选为 A20。全体包括高海拔子集，不能把两行视作独立重复。阈值的命中和误报见\hyperlink{diagnosis}{第5节}。\src{P8}}
\textbf{读图动作。} 问听众为何 A20 在固定阈值命中更多的同时，AP20 仍下降。如果只给出 AP20 图，请补分母、固定阈值计数及原先接受的误报线。\par
\textbf{补齐原筛查。} A20 两项通过的是 A 的命中与真强雨 q32 越线。强雨物理 Pinball/20 从 0.506865 升至 0.628274，非强雨 q32 越线从 5,769 增至 352,680，弱雨损失和发生指标也退化。\src{P13}\par



\clearpage
\hypertarget{freeze}{}
\section*{把门槛写成可执行的协议}
\label{p-freeze}
{\color{accent}\bfseries 2026-11-21 的 IMERG 开发评价冻结实例}\par

\textbf{比较器与身份。} B7/REAL/COSINE9，seed2029、固定 E9。共同有效监督像元时 3,303,090，参考有雨 407,820，严格大于 20 的强雨 717，弱雨 \(0.1<R\le1\,\mmh\) 共有 246,473。先核对真值、掩膜、海拔分区和概率节点合同，再算差值。\src{P10}

\begin{tabularx}{\linewidth}{@{}L L@{}}
\toprule
事前要求 & 它约束的失败模式 \\
\midrule
强雨物理 \term{pinball}{Pinball/20} 严格下降 & 防止只把少量输出推过阈值，却没有改善条件幅值。\\
真强雨 q32 越线 \(\ge160\)，非强雨 q32 越线 \(\le5{,}769\) & 防止单纯整体提高上尾，同时放大非强雨位置的尾部支撑。\\
固定 A 的命中 \(>0\)，误报 \(=0\) & 要求固定读出恢复检出，同时守住当前冻结的误报线。\\
弱雨 log Pinball、发生 Brier 不升，发生 AP、AUROC 不降 & 检查候选没有通过牺牲较常见降水与发生表现取得局部改善。\\
\bottomrule
\end{tabularx}

\textbf{实际算分顺序。} 对强雨真值样本，物理 Pinball/20 先在 32 个物理分位上平均，再在强雨像元时上平均。固定 \term{readout}{A 点值}是发生概率乘以 32 个物理分位数的算术平均，以 64 位浮点计算，转为 32 位浮点后按严格 \(A>20\) 判定。AP20 用同一 A 分数，是全域和高海拔都需报告的排序信息，在本次新协议中属于辅助报告。历史 A20 按该实验的预设排序停止规则封存。\src{P8}\src{P10}

\boxnote{失败也要留下什么}{记录尝试改了什么、为何可能有效、哪项结果推翻预期、哪些结果仍成立，以及下一项需要什么证据。当前冻结了开发筛查和最终测试边界，待核实连续事件覆盖；截至本快照尚无准入的新候选，2025 的实际测试范围与执行仍待研究者绑定。\src{P10}}

\sub{把实验阶段看作逐项排除解释}
B0/B1 比较时序，B2/B3 比较多光谱输入条件，B3/B4 检查 GFS 环境增量，随后 B5/B6 引入静态地形，B7/B8 探索动态地形与融合方式。各段有自己的比较器、人口和种子。列出这条路线时，逐项报告实际匹配的比较和各自反例。\src{P4}\src{P11}


\chapterpage{verification}{从像元到暴雨过程}{先确认参考口径，再扩大验证范围}
\label{p-verification}
\sub{同源参考和时间留出分别回答什么？}
当前采用 \term{reference}{IMERG V07 Final 同源参考}。它也提供学习标签，因此配对结果衡量对该参考的表现。最终候选和协议冻结后，2025 可用于同源时间留出检验；截至案例日期仍封存。IMERG 的 V07 是版本，Final Run 是处理流，2025 Final Test 是数据划分。\src{P1}\src{R1}

\sub{为什么 717 个像元时没有给出 717 场暴雨？}
像元时是一个目标窗内的一个像元；场景是一次格点场采样；天气过程跨多个相邻时刻和空间位置。2024 Dev 的 717 个强雨像元时分布于 207 个含强雨窗口，多数相邻场景约隔 6 小时。仅凭稀疏窗口无法恢复完整暴雨过程，也不能把相邻像元视作独立样本。\src{P1}

事件评价要先用连续参考建立\term{event}{过程目录}，事前写清邻域、时间槽、位移容忍、最小面积、持续时间、缺测和截断。再映射共同有效的模型窗，分别报告目录过程数、可匹配数、可评价数及匹配覆盖。当前协议为草案，覆盖待核实。若仅能核对稀疏时窗，就把结论限定在匹配窗内。\src{P1}

\sub{什么能增强独立物理验证？}
站点给出点测，雷达 \term{qpe}{QPE} 是经过估计的面积结果，\term{dpr}{DPR} 是卫星过境时的瞬时足迹。纳入前检查来源链与标签是否共享观测、时间支撑、空间代表性、地形遮挡、版本和缺测。现行路线已确认 IMERG 参考并冻结开发评价，站点、雷达和 DPR 为可选补充，独立观测证据仍为 \texttt{PARTIAL}。\src{P1}\src{P10}\src{R1}

\sub{怎样报告不确定性}
训练随机性用同规则的配对 seed 差值描述。天气差异用可核对的过程或时块来表述。若进行重采样，两臂在同一抽样中的事件集合要一致。报告覆盖、样本数、差值区间和正负方向。条件不满足时保留描述性效应，说明哪项不确定性仍未估计。

\boxnote{检验一条结论的适用范围}{问自己「换年份、换天气过程、换观测来源」三件事是否都实际做过。每多写一个适用范围，都要指出相应比较的样本、协议与结果。}

\chapterpage{communication}{写结果、讲组会、回答追问}{把研究结论写成别人能逐项核查的句子}
\label{p-communication}
\sub{可直接改写成结果段的案例}
「在 2024 Dev 的 963 个配对场景、3,303,090 个共同有效像元时中，比较固定 E9 的 COSINE9 与 A20 点值。A20 在严格 \(>20\,\mmh\) 的阈值下命中 627/717 个强雨像元时，并产生 177,098 次误报。全体 AP20 从 0.018691 降至 0.007786，高海拔子集同向下降，因此按事前准则封存该候选。这一比较衡量所用开发场景对 IMERG 参考的点值检出和排序。」\src{P1}\src{P8}

\sub{三分钟组会分配}
\begin{tabularx}{\linewidth}{@{}p{2.15cm}L@{}}
\toprule
用时 & 要讲出的句子 \\
\midrule
0:00--0:35 & 我们想估计云南山区半小时降水。强雨点值长期难以过阈值，因此这一轮要分辨幅值不足和越线后误报。\\
0:35--1:15 & 两臂使用相同场景、真值与掩膜，我用固定该阈值的命中和误报检验可用性，再用 AP20 看强雨在分数中的排序。\\
1:15--2:20 & 报告参考臂、候选、分母和误报。指出 AP20 下降与 AUROC20 小幅上升可以同时发生，并说明固定阈值验收为何失败。\\
2:20--3:00 & 这仍是对 2024 Dev 上同源参考的比较。下一份能够改变结论的材料是事前定义的、覆盖可核验的过程评价或留出结果。\\
\bottomrule
\end{tabularx}

\sub{给听众留一份复核记录}
在图注或一页说明上，保留\term{claim}{主张与证据}、输入和参考版本、时空窗、两臂唯一差异、阈值、分母、效应方向、反例与局限。附可定位的公开收据。写“改善”时附绝对数和变化方向，写“显著”前给出推断方法、独立性假设及区间。

\sub{一页说明写什么}
\textbf{标题}用一个问题。左侧写对照身份、实际改变和评价范围；右侧放带分母的结果，至少保留一项反例。末行写当前结论、可推翻它的下一项观测，以及供听众查证的收据编号。\back{sources}{证据索引}

\boxnote{组会常见追问}{「为什么不用 2025 选个更好的阈值？」开发选择必须和最后一次检验隔离。「为什么有雨概率高还漏强雨？」发生与条件强度是不同输出，固定 A 还取决于条件均值。「为什么现有结果不能证明地形机制？」因为目前局部变化、反例与广泛复现的证据范围不同。}

\chapterpage{exercises}{练习与自评}{闭卷完成，再翻到参考答案}
\label{p-exercises}
\sub{练习一　设计一项可被否证的研究}
选一类环境字段或一个新增卫星通道，写 150 字问题卡。交付比较臂、样本身份、主指标、至少两条保护线、失败条件，以及失败之后需要核对的一个原因。然后把同一模板迁移到你的课题。

\sub{练习二　审查时间和单位}
教学例：目标窗为 12:00--12:30 UTC。一场记录累计 11 mm；某个有效时刻为 12:30 的模式产品在 12:42 才可下载，而业务拟在 12:35 发布。写出平均雨强、是否为强雨，以及该资料能否用于这次业务发布。说明时窗完整性有何影响。

\sub{练习三　算出研究结论的代价}
用案例中的 717 个强雨像元时、A20 的 627 命中和 177,098 误报，求漏报数、POD 和精确率。对照全体 AP20，写一句 60 字以内的结果，带上数据范围和停止理由。

\sub{练习四　把一段过强的结果写回证据}
改写「加入动态地形后，模型能准确检测山区所有暴雨，已可用于预警。」必须引用海拔层差异、像元时与事件的区别、固定点值结果，并留下业务可用性条件。随后找同学当评委，完成三分钟讲解。

\bigskip
\hyperlink{answers}{核对参考答案 $\rightarrow$}

\clearpage
\hypertarget{answers}{}
\section*{参考答案与输出评分}
\textbf{一。} 示例：在固定共同开发样本和固定检查点比较「原环境输入」与「只新增一个水汽字段」。以参考强雨物理 Pinball 为主指标，同时保护弱雨分位损失与固定阈值误报。若主指标无改善或任一保护线越界，结束该候选，优先查单位和地下层掩膜。具体容差需在看结果前确定。

\textbf{二。} 11 mm 除以 0.5 h 等于 22 mm/h，满足严格 \(>20\)。产品 12:42 才到达，晚于 12:35，不符合这次发布。若累计窗覆盖不完整，标记缺测并核对资格，不能按半小时直接换算。

\textbf{三。} 漏报 \(717-627=90\)，POD \(=627/717\approx87.45\%\)，精确率 \(=627/(627+177098)\approx0.35\%\)。示例：「2024 开发配对中，A20 虽命中 627/717，误报 177,098 且 AP20 下降，未通过冻结准则。」对照臂无正预测时，精确率的分母为零，报告未定义或注明工具的约定。

\textbf{四。} 示例：「2024 Dev 的单 seed 配对中，REAL 在高海拔强雨像元时的物理分位损失略降，两个较低海拔层退化。两臂固定点估计在强雨阈值下均零命中，事件覆盖和业务可用性仍待核实。」\src{P6}

\sub{同伴评分表}
\begin{tabularx}{\linewidth}{@{}p{3.4cm}L@{}}
\toprule
维度 & 满分 2 分的可观察表现 \\
\midrule
问题与设计 & 能给出可被推翻的主张、匹配对照和事前保护线。 \\
数据与时空 & 能说清参考、目标窗、单位、共同样本和资料可用时间。 \\
读数与解释 & 能复算分母，区分分位损失、点值检出、概率和排序。\\
表达与追溯 & 能讲清结论边界和反例，并指出两份可核对证据。\\
\bottomrule
\end{tabularx}
每项可独立完成记 2 分，提示后完成记 1 分，无法完成记 0 分。教学建议是 7/8 分且数据时空、读数两项均满分后，再尝试独立讲解。选一条自己的研究记录，交出一页问题卡、一段结果和一次三分钟讲解作为迁移测验。
